% !PRINTED 523
\weberpage{523}

adjoin, in which
\weq{5}{D = + a^2b^2 + 18\,abc - 4a^3c - 4b^3 - 27c^2}
is the discriminant of the equation (3), then the group reduces
to (2).

In the field $\Omega'$, which arises from $\Omega$ by this
adjunction, every root can be expressed rationally in terms of every
other, for we have
\[
\begin{gathered}
\alpha_1 + \alpha_2 = a - \alpha\\
\alpha_1 - \alpha_2 = \dfrac{\sqrt{D}}{f'(\alpha)},
\end{gathered}
\]
hence
\weq{6}{\begin{gathered}
2\alpha_1 = a - \alpha + \dfrac{\sqrt{D}}{f'(\alpha)}\\
2\alpha_2 = a - \alpha - \dfrac{\sqrt{D}}{f'(\alpha)},
\end{gathered}}
and in these the permutations $\pi$, $\pi^2$ can be carried out.
The cubic equation is therefore, after adjunction of $\sqrt{D}$,
its own \gesperrt{Galois} resolvent.

All this remains valid if the domain of rationality is any
special field $\Omega$ in which $a$, $b$, $c$ and
$\sqrt{D}$ are contained, unless $f(x)$ is itself reducible in
$\Omega$, and so has a rational root. For besides itself
and the cyclic group (2), the group (1) has only
intransitive divisors, namely the unit group and three groups
of the type (1, 2).

If we now wish to solve the cubic equation, that is, to reduce it
to a pure cubic equation, we must seek a
function $v$ of the roots which does not itself admit the cyclic
permutation $\pi$, but whose cube does. Such a function
can exist, however, only if the third roots of unity, or, what
is the same thing, $\sqrt{-3}$, is adjoined to the field $\Omega$; for
$v$ must, on application of $\pi$, acquire a third root of unity as
a factor.

By this adjunction no reduction of the group can
occur, because the group (2) has no divisor besides the
unit group, and the reduction to the unit group can
take place not through a quadratic, but only through a cubic
equation (\S.\,157). We therefore set
\[
\varepsilon = \dfrac{-1+\sqrt{-3}}{2}, \qquad
\varepsilon^2 = \dfrac{-1-\sqrt{-3}}{3}
\]
% [sic] The source prints 3 in the denominator of \varepsilon^2; algebra requires 2. The printed value is preserved source-faithfully.
%CONT>>
